\specialchapter{参考文献}

\begin{enumerate}
\def\labelenumi{(\Roman{enumi})}
\item
  J. HESSEL, Krysallometrie oder Krystallonomie und Krystallographie（1830，重印于 Ostwald's Klassiker，第 88 和 89 号，Leipzig（Teubner），1897）
\item
  A. Bravais，《关于对称形多面体的回忆录》，Journal de Math.，（1），v. XIV（1849），p.~141-180。
\item
  A. MÖBIUS: a) Ueber das Gesetz der Symmetrie der Krystalle und die Anwendung dieses Gesetze auf die Eintheilung der Krystalle in Systeme, J. für die reine und angewandte Math., v. XLIII (1852), p.~365-374 (= Gesammelte Werke, v. II, Leipzig (Hirzel), 1886, p.~349-360); b) Theorie der symmetrischen Figuren, Gesammelte Werke, v. II, Leipzig (Hirzel), 1886, p.~561-708.
\item
  L. SCHLÄFLI, Theorie der vielfachen Kontinuität. Denkschr. der Schweiz. naturforsch. Gesellschaft, v. XXXVIII (1901), p.~1-237 (= Ges. math. Abhandlungen, v. I, Basel (Birkhäuser), 1950, p.~167-387).
\item
  W. R. HAMILTON，《关于单位根新系统的备忘录》，Phil. Mag.，（4），v. XII（1856），p.~446。
\item
  C. JORDAN，《关于运动群的回忆录》，Ann. di Mat.，v. Il（1868-69），p.~167-215 和 322-345（= Oeuvres，v. IV，Paris（Gauthier-Villars），1964，p.~231-302）。
\item
  E. GOURSAT，《关于正交代换与空间的正则划分》，Ann. Ec. Norm. Sup.，（3），v. VI（1889），p.~9-102。
\item
  W. KILLING, Die Zusammensetzung der stetigen endlichen Transformationsgruppen: I) Math. Ann., v. XXXI (1888), p.~252-290; II) ibid., v. XXXIII (1889), p.~1-48; III) ibid., v. XXXIV (1889), p.~57-122; IV) ibid., v. XXXVI (1890), p.~161-189.
\item
  E. CARTAN：a）《关于有限连续变换群的结构》（学位论文）。Paris（Nony），1894（= Oeuvres complètes，Paris（Gauthier-Villars），1952，v. I₁，p.~137-287）；b）《关于将有限连续变换群的结构约化为其规范形式》，Amer. Journ. of Math.，v. XVIII（1896），p.~1-46（= Oeuvres complètes，Paris（Gauthier-Villars），1952，v. I₁，p.~293-353）；c）《对偶原理与单群和半单群理论》，Bull. Sci. Math.，v. XLIX（1925），p.~361-374（= Oeuvres complètes，Paris（Gauthier-Villars），1952，v. I₁，p.~555-568）；d）《单群的几何学》，Ann. di Mat.，（4），v. IV（1927），p.~209-256（= Oeuvres complètes，Paris（Gauthier-Villars），1952，v. I₂，p.~793-840）；e）《具有单基本群的几何的某些显著黎曼形式》，Ann. Ec. Norm. Sup.，（3），v. XLIV（1927），p.~345-467（= Oeuvres complètes，Paris（Gauthier-Villars），1952，v. I₂，p.~867-989）；f）《对回忆录 ``单群的几何学'' 的补充》，Ann. di Mat.，v. V（1928），p.
\end{enumerate}

253-260 (= Oeuvres complètes, Paris (Gauthier-Villars), 1952, v. I₂, p.~1003-1010).

\begin{enumerate}
\def\labelenumi{(\Alph{enumi})}
\setcounter{enumi}{23}
\tightlist
\item
  R. FRICKE 和 F. KLEIN，Theorie der automorphen Funktionen，Leipzig（Teubner），1897。
\end{enumerate}

\begin{enumerate}
\def\labelenumi{(\Roman{enumi})}
\setcounter{enumi}{10}
\item
  L. E. DICKSON：a）《任意域上的线性群理论》，Trans. Amer. Math. Soc.，v. II（1901），p.~363-394；b）《单群的新系统》，Math. Ann.，t. LX（1905），p.~137-150；c）《任意域中与三次曲面上 27 条直线构形有关的一类群》，Quart. Journ. of Math.，v. XXXIII（1901），p.~145-173 和 v. XXXIX（1908），p.~205-209。
\item
  A. SPEISER. Theorie der Gruppen von endlicher Ordnung, Berlin (Springer), 1924.
\item
  H. WEYL, Theorie der Darstellung kontinuerlicher halb-einfacher Gruppen durch lineare Transformationen, Math. Zeitschr., v. XXIII (1925), p.~271-309, v. XXIV (1926), p.~328-395 and 789-791 (= Selecta, Basel-Stuttgart (Birkhäuser), 1956, p.~262-366).
\item
  H. S. M. COXETER：a）《基本域为单纯形的群》，Journ. Lond. Math. Soc.，v. VI（1931），p.~132-136；b）《具有正棱柱图形的多胞体》，Proc. Lond. Math. Soc.，（2），v. XXXIV（1932），p.~126-189；c）《由反射生成的离散群》，Ann. of Math.，（2），v. XXXV（1934），p.~588-621；d）《形如 \(\mathbf{R}_i^2 = (\mathbf{R}_i\cdot \mathbf{R}_j)^{k_{ij}} = 1\) 的有限群的完全枚举》，Journ. Lond. Math. Soc.，v. X（1935），p.~21-25；e）《正多胞体》，New York（Macmillan），1948（第 2 版，1963）；f）《由反射生成的有限群的生成元之积》，Duke Math. Journ.，v. XVIII（1951），p.~765-782。
\end{enumerate}

（XIV bis）H. S. M. COXETER 和 H. WEYL，《连续群的结构与表示》（Inst. for Adv. Study，由 N. Jacobson 和 R. Brauer 编印的讲义，1934-35）：附录。

\begin{enumerate}
\def\labelenumi{(\Roman{enumi})}
\setcounter{enumi}{14}
\item
  H. S. M. COXETER 和 W. O. J. MOSER，《离散群的生成元与关系》，Ergebn. der Math.，Neue Folge，Bd. 14，Berlin（Springer），1957（第 2 版，1963）。
\item
  B. L. van der WAERDEN, Die Klassification der einfachen Lieschen Gruppen, Math. Zeitschr., v. XXXVII (1933). p.~446-462.
\item
  E. WITT, Spiegelungsgruppen und Aufzählung halbeinfacher Lieschen Ringe, Abh. Math. Sem. Hamb. Univ., v. XIV (1941), p.~289-322.
\item
  E. STIFEL, Ueber eine Beziehung zwischen geschlossenen Lie'sche Gruppen und diskontinuierlichen Bewegungsgruppen euklidischer Räume und ihre Anwendung auf die Aufzählung der einfachen Lie'schen Gruppen, Comm. Math. Helv., v. XIV (1941-42), p.~350-380.
\item
  C. CHEVALLEY：a）《单李代数及其表示的分类》，C. R. Acad. Sci.，v. CCXXVII（1948），p.~1136-1138；b）《由反射生成的有限群的不变量》，Amer. Journ. of Math.，v. LXXVII（1955），p.~778-782；c）《关于某些单群》，Tohôku Math. Journ.，（2），v. VII（1955），p.~14-66；d）《代数李群的分类》，2 卷，Paris（Inst. H. Poincaré），1956-58。
\item
  HARISH CHANDRA：a）《半单李代数的普遍包络代数的一些应用》，Trans. Amer. Math. Soc.，v. LXX（1951），p.~28-96；b）《关于 Bruhat 的一个引理》，Journ. de Math.，（9），v. XXXV（1956），p.~203-210。
\item
  F. BRUHAT，《复半单李群的诱导表示》，C. R. Acad. Sci.，v. CCXXXVIII（1954），p.~437-439。
\item
  A. BOREL，《线性代数群》，Ann. of Math.，（2），v. LXIV（1956）。p.~20-80。
\item
  A. J. COLEMAN，《单群的 Betti 数》，Can. Journ. of Math.. v. X（1958）。p.~349-356。
\item
  R. STEINBERG。《有限反射群》。Trans. Amer. Math. Soc.，v. XCI（1959）。p.~493-504。
\item
  J. TITS：a）《单群及其相关几何》，Proc. Int. Congress Math.，Stockholm，1962，p.~197-221；b）《Bruhat 定理与抛物子群》。C. R. Acad. Sci.，v. CCLIV（1962），p.~2910-2912；c）《代数单群与抽象单群》。Ann. of Math.，（2），v. LXXX（1964），p.~313-329。
\item
  N. IWAHORI 和 H. MATSUMOTO，《关于 Bruhat 分解以及 \(p\) 进 Chevalley 群 Hecke 环的结构》，Publ. Math. I.H.E.S.，no. 25（1965），p.~5-48。
\item
  A. BOREL 和 J. TITS。《约化群》。Publ. Math. I.H.E.S.，no. 27（1965），p.~55-150。
\item
  F. BRUHAT 和 J. TITS。《局部域上的单代数群》，Proc. Conf. on local Fields，p.~23-36。Berlin（Springer），1967。
\item
  R. CARTER，《单群与单李代数》，Journ. Lond. Math. Soc.，v. XL（1965），p.~193-240。
\end{enumerate}

